\section{Introduction} %NOT DONE
%Describe OSPF
Telecommunication backbone networks consist of groups of nodes within the same domain that carry aggregated traffic across geographic regions, interconnecting cities and networks through high-capacity links. The task of a telecommunication backbone network is to transmit data between sources and destinations. These networks typically operate under dynamic topologies managed by telecommunications companies and are referred to as intra-domain networks. Intra-domain networks require routing protocols to facilitate packet transmission between nodes. A widely used protocol in these networks is the Open Shortest Path First (OSPF) routing protocol \cite{rfc2328}. OSPF is a link-state protocol that detects topological changes within the Autonomous System (AS), routes packets based on the IP header, and determines the optimal path using Dijkstra's algorithm.

%Describe OSPF Limitation
Although OSPF is capable of routing packets, it has limitations in data delivery. The authoritative OSPF standard employs a single configuration cost, which restricts protocol adjustments to link weights only \cite{rfc2328}. OSPF does not account for the real-time condition of links between nodes. When static weights are used, the protocol ignores real-time load and capacity, which can result in network congestion \cite{1039866}. Consequently, these limitations reduce OSPF's flexibility for intra-domain routing. Telecommunication backbones face challenges in maintaining connectivity because standard OSPF offers only limited traffic-engineering capability, confined to link-weight tuning.


%Describe Research Gap  
As the network traffic demand grows, there is a growing need to find more adaptive solutions for network routing. Multi-Agent Reinforcement Learning with a Graph Neural Network (GNN) method has recently been applied to dynamic routing in mobile ad-hoc networks (MANETs) \cite{11134361}. However, existing research primarily addresses wireless, mobile, and topology-varying environments. Their applicability to fixed intra-domain telecommunication backbones where routing is capacity-constrained and congestion-driven remains unexplored. Furthermore, prior studies have primarily evaluated routing performance using synthetic, mobility-driven traffic rather than real network demand.

%Describe Proposed Approach
To address these gaps, this work applies and adapts Multi-Agent Reinforcement Learning with a Graph Neural Network (GNN) backbone to intra-domain telecommunication backbone routing. Reinforcement Learning is a machine learning paradigm in which agents learn optimal actions through a system of rewards and punishment \cite{SHAKARAMI2020107496}. Multi-agent Reinforcement Learning involves a set of agents interacting with each other in the environment, each doing trial and error to maximize reward \cite{marl-book}. Combining with the GNNs, which are designed to process and analyze graph structure data \cite{10960451}, the network topology can be effectively represented as a graph.

%Describe Contribution
This study evaluates the performance of Multi-Agent Reinforcement Learning (MARL) with Graph Neural Network (GNN) integration against Open Shortest Path First (OSPF). The evaluation is conducted on three distinct telecommunication backbone topologies to enable a comprehensive comparison. Topology and traffic data are sourced from the Survivable Network Design Library \cite{SNDlib10} \cite{OrlowskiPioroTomaszewskiWessaely2010}. Regarding MARL training and execution, this work adopts a centralized training and decentralized execution framework. The work employs analytical surrogate methods for training due to resource constraints and performs packet-level evaluations using NS-3 \cite{Riley2010}. Furthermore, this work measures performance by collecting Maximum Link Utilization, Packet Loss, and Delay as metrics for comparison with OSPF.  

The main contributions of this work are the following:
\begin{enumerate}
\item The application and adaptation of Multi Agent Reinforcement Learning with Graph Neural Network for telecommunication backbone networks. 
\item A diagnosis of the conditions that cause MARL to fail and degrade routing performance.
\item A reproducible real-data evaluation pipeline of the multi-agent reinforcement learning algorithm using ns-3 simulator.
\end{enumerate}

